\section{Continuous equivalent layers}
\label{sec:continuous_eqlayer}

We now move from a discrete to a continuous formulation of the equivalent layer with 
the purpose of discussing its theoretical properties.
Let $d^{o}(\mathbf{r})$ be a harmonic function representing the observed potential field 
produced by the 3D sources at a generic observation point 
$\mathbf{r} = ( x , y , z )^{\top} \notin \overline{\mathcal{V}}$ located above them.
Besides, consider that $d^{o}(\mathbf{r})$ can be approximated within a given accuracy $\varepsilon$ 
by the potential field due to a continuous and planar physical-property 
distribution $p(\mathbf{r}'')$ defined at a horizontal plane $z''_{0}$ located below the 
observations (Figure \ref{fig:geometry}). 
As in the discrete case, it is important bear in mind that the plane $z''_{0}$ may cross the sources.
From these definitions, we propose the following definition:

\begin{definition}[\textbf{Continuous equivalent layer}]
	Let us assume to have an observed potential field $d^{o}(\mathbf{r})$ defined above 
	a horizontal surface $\mathcal{S}''$ at $z''_{0}$, formed by points 
	$\mathbf{r}'' = ( x'' , y'' , z''_{0} )^{\top}$. 
	Besides, consider (i) a horizontal surface $\mathcal{S}^{\dagger}$ formed by points 
	$\mathbf{r}^{\dagger} = ( x^{\dagger} , y^{\dagger} , z^{\dagger}_{0} )^{\top}$ at a constant 
	$z^{\dagger}_{0}$ satisfying $z^{\dagger}_{0} < z''_{0}$ and $z^{\dagger}_{0} < z^{t}_{0}$, 
	with $z^{t}_{0}$ defining the depth to the sources top, and (ii) a harmonic kernel 
	$a(\mathbf{r}, \mathbf{r}'')$ that depends only on $\mathbf{r}$ and $\mathbf{r}''$, as well as 
	vanishes as the distance $\| \mathbf{r} - \mathbf{r}'' \|_{2}$ increases. 
	Then, for a chosen accuracy $\varepsilon > 0$, we may call $p_{\varepsilon}(\mathbf{r}'')$ 
	a \textbf{continuous equivalent layer} if its predicted potential field 
	\begin{equation}
		d_{\varepsilon}(\mathbf{r})
		= \int\limits_{-\infty}^{\infty}\int\limits_{-\infty}^{\infty}
		p_{\varepsilon}(\mathbf{r}'') \, a(\mathbf{r}, \mathbf{r}'') \:\: d\mathcal{S}''
		\: , \quad z < z''_{0} \: ,
		\label{eq:predicted-data-integral}
	\end{equation}
	satisfies 
	\begin{equation}
		\int\limits_{-\infty}^{\infty} \int\limits_{-\infty}^{\infty}
		\vert d^{o}(\mathbf{r}^{\dagger}) - 
		d_{\varepsilon}(\mathbf{r}^{\dagger}) \vert^{2}
		\:\: d\mathcal{S}^{\dagger} \le \varepsilon \: ,
		\quad z^{\dagger}_{0} < z^{t}_{0} \: , \quad z^{\dagger}_{0} < z''_{0} \: .
		\label{eq:epsilon-bound-data}
		\end{equation}
	\label{def:continuous-eqlayer}
\end{definition}

In Definition \ref{def:continuous-eqlayer}, $p_{\varepsilon}(\mathbf{r}'')$ defines a continuous physical-property distribution 
that extends to infinite along the horizontal directions and the kernel function is usually defined by
$a(\mathbf{r}, \mathbf{r}'') = D_{a} \Psi(\mathbf{r}, \mathbf{r}'')$, as a generic partial derivative of 
the function $\Psi(\mathbf{r}, \mathbf{r}'')$ (equation \ref{eq:Psi}),
where $D_{a}$ is defined in terms of the space-domain differential operator 
$D$ (equation \ref{eq:diff-operators-generic}).

Similarly to the discrete case (equation \ref{eq:transformed-field-layer-space}), a particular continuous equivalent layer 
$\tilde{p}_{\varepsilon}(\mathbf{r}'')$ leads to a transformed potential field given by:
\begin{equation}
	f_{\varepsilon}(\mathbf{r}^{\ddagger})
	= \int\limits_{-\infty}^{\infty}\int\limits_{-\infty}^{\infty}
	\tilde{p}_{\varepsilon}(\mathbf{r}'') \, t(\mathbf{r}^{\ddagger}, \mathbf{r}'') \:\: d\mathcal{S}''
	\: , \quad z < z''_{0} \: ,
	\label{eq:transformed-field-integral-space}
\end{equation}
where $\mathbf{r}^{\ddagger} = (x^{\ddagger}, y^{\ddagger}, z^{\ddagger})^{\top} \notin \overline{\mathcal{V}}$ 
are points located above the layer plane $z''_{0}$, but that can eventually 
coincide with those containing the observed potential field; 
$f_{\varepsilon}(\mathbf{r}^{\ddagger})$ approximates 
$f(\mathbf{r}^{\ddagger})$ (equations \ref{eq:transformed-field-space}); and
$t(\mathbf{r}^{\ddagger}, \mathbf{r}'') = D_{t} \Psi(\mathbf{r}^{\ddagger}, \mathbf{r}'')$ is the transformation kernel,
which is usually defined as a generic partial derivative of the potential function $\Psi$
(equation \ref{eq:Psi}), with $D_{t}$ given in terms of $D$ (equation \ref{eq:diff-operators-generic}).